\SourceLeaf{292}%
Länge einer Dreiecksseite \(r\), ihrem Azimut \(T\) am Anfangspunkt und dessen Breite \(S\) sind das Azimut \(T'\pm180^0\) der Seite am anderen Endpunkt, dessen Breite \(S'\) und der Längenunterschied \(\lambda\) der beiden Punkte abzuleiten. Dabei geht es lediglich um die Auflösung eines sphärischen Dreiecks. Die Aufgabe verdient hier nur deshalb einen Platz, weil die üblichen Formeln etwas umgeformt werden müssen, wenn man nach der Bemerkung in Artikel 10 für die Ergebnisse dieselbe Genauigkeit erreichen will, mit der \(r\) gegeben ist, ohne Logarithmen mit mehr Stellen heranzuziehen. Damit man je nach Bedarf zwischen den verschiedenen Lösungsarten wählen kann, gebe ich zuerst diejenigen an, die auf den bekannten elementaren Formeln der sphärischen Trigonometrie beruhen.

Erste Methode
\[
\begin{aligned}
\operatorname{tang}s&=\cos T\operatorname{tang}r\\
\operatorname{tang}\lambda&=\frac{\operatorname{tang}T\sin s}{\cos(S-s)}\\
\operatorname{tang}S'&=\cos\lambda\operatorname{tang}(S-s)\\
\sin T'&=\frac{\sin T\cos S}{\cos S'}
\end{aligned}
\]
\indent Zweite Methode
\[
\begin{aligned}
\operatorname{tang}R&=\frac{\operatorname{tang}S}{\cos T}\\
\operatorname{tang}T'&=\frac{\operatorname{tang}T\cos R}{\cos(R-r)}\\
\operatorname{tang}S'&=\cos T'\operatorname{tang}(R-r)\\
\sin\lambda&=\frac{\sin r\sin T}{\cos S'}=\frac{\sin r\sin T'}{\cos S}
\end{aligned}
\]
\indent Dritte Methode
\[
\begin{aligned}
\sin(45^0+\tfrac{1}{2}S')\sin\tfrac{1}{2}(T'+\lambda)&=\sin(45^0+\tfrac{1}{2}(S+r))\sin\tfrac{1}{2}T\\
\sin(45^0+\tfrac{1}{2}S')\cos\tfrac{1}{2}(T'+\lambda)&=\sin(45^0+\tfrac{1}{2}(S-r))\cos\tfrac{1}{2}T\\
\cos(45^0+\tfrac{1}{2}S')\sin\tfrac{1}{2}(T'-\lambda)&=\cos(45^0+\tfrac{1}{2}(S+r))\sin\tfrac{1}{2}T\\
\cos(45^0+\tfrac{1}{2}S')\cos\tfrac{1}{2}(T'-\lambda)&=\cos(45^0+\tfrac{1}{2}(S-r))\cos\tfrac{1}{2}T
\end{aligned}
\]
\indent Was die Kürze der Rechnung angeht, hat die dritte Methode einen gewissen Vorteil gegenüber den beiden anderen, während diese im Allgemeinen etwas genauere Ergebnisse liefern können, insbesondere für \(\lambda\), das sich damit stets mit völlig ausreichender Genauigkeit bestimmen lässt. Dagegen wird \(T'\) durch die erste Methode vergleichsweise ungenau bestimmt, wenn es nahe an einem rechten Winkel liegt. Verlangt man aber alle drei Ergebnisse mit gleichmäßiger und aus der Sicht von Artikel 10 ausreichender Genauigkeit,
